\documentclass[11pt]{amsart}
\usepackage[T1]{fontenc}
\usepackage[utf8]{inputenc}
\usepackage{amssymb}
\usepackage[margin=1.05in]{geometry}
\usepackage[hidelinks]{hyperref}
\usepackage{longtable}
\newcommand{\stamptime}{\the\numexpr\time/60\relax:\ifnum\numexpr\time-60*(\time/60)\relax<10 0\fi\the\numexpr\time-60*(\time/60)\relax}
\AddToHook{shipout/foreground}{\put(0.5\paperwidth,-\paperheight+0.55cm){\makebox[0pt]{\tiny\ttfamily unit/doc/\jobname.tex \ --- \ compiled \today, \stamptime}}}
\newtheorem{theorem}{Theorem}
\newtheorem{proposition}[theorem]{Proposition}
\newtheorem{lemma}[theorem]{Lemma}
\newtheorem{corollary}[theorem]{Corollary}
\theoremstyle{definition}
\newtheorem{definition}[theorem]{Definition}
\theoremstyle{remark}
\newtheorem{remark}[theorem]{Remark}
\newcommand{\Om}{\Omega}
\newcommand{\code}[1]{\texttt{#1}}
\newcommand{\cT}{\mathsf{isTrunc}}
\newcommand{\cC}{\mathsf{isConn}}
\newcommand{\refl}{\mathsf{refl}}
\newcommand{\ap}{\mathsf{ap}}
\newcommand{\el}{\mathsf{el}}
\newcommand{\Fib}{\mathrm{Fib}}
\newcommand{\fib}{\mathrm{fib}}
\newcommand{\tr}[1]{\|#1\|}
\newcommand{\Ab}{\mathrm{Ab}}
\newcommand{\St}{\mathsf{Stage}}
\newcommand{\Tor}{\mathsf{Tor}}
\newcommand{\SC}{\mathsf{SC3}}
\newcommand{\Tri}{\mathsf{Triples}}
\newcommand{\Fam}{\mathsf{Fam}}
\newcommand{\Mu}{\mathsf{Mu}}
\newcommand{\pteq}{\simeq_*}
\newcommand{\U}{\mathcal{U}}
\title{Pointed simply connected 3-types are classified by triples $(A,B,k)$}
\date{28 October 2026}
\begin{document}
\maketitle

\noindent\textbf{Purpose.} Informal proof for
\code{doc/postnikov-classification-handoff.txt}: an equivalence of types
\[ \SC_\ell\ \simeq\ \Tri_\ell \]
whose inverse is the realisation $(A,B,k)\mapsto(\Fib k,\ldots)$ of
\code{PostnikovRealization.Triple}. The cocycle $k$ is an actual pointed map;
nothing is truncated or quotiented.

\noindent\textbf{Attribution.} The proof is the equivalence between
$(n{+}1)$-groups and untruncated Sính triples of Perry Hart and Owen Milner,
\emph{Classifying 2-Groups in Homotopy Type Theory}, LICS 2026, LIPIcs vol.~380,
DOI \code{10.4230/LIPIcs.LICS.2026.55}, \S\S8--9 (Lemma 40, Theorem 41, Note 42),
specialised to $n=2$ and to simply connected types. The paper states Theorem 41 as a
result of the paper; the repository's \code{README.md} and \code{Sinh/README.md}
call it ``Owen Milner's type equivalence'', and we follow that attribution. The
delooping and Eilenberg--MacLane machinery of Unit (\code{sag/em}) is David
Wärn's (DW1); the sphere universal-property strategy of the existing
development is Wärn's, and plays no role here.

\section{Sources read}

\noindent\textbf{Paper:} \code{CANONICITY/2groups.pdf} (25 pp., LICS 2026 version),
\S8 (Definition 35, Theorem 36, Example 37, Theorem 38, Corollary 39) and \S9
(Lemma 40, Theorem 41, Note 42, Theorem 43).

\noindent\textbf{Agda:} \code{https://github.com/PHart3/2-groups-agda}, branch
\code{lics}, commit \code{9ec62dd454a8330278fdfce5cac9d1a2f3cc48de} (9 July
2026), cloned read-only (shallow) into the session scratchpad; not built. Files
read in full: \code{Sinh/Sinh-classif.agda}, \code{Sinh/Sinh-action.agda},
\code{Sinh/Sinh-cocycle.agda}, \code{Final/Final-thms.agda},
\code{HoTT-Agda/core/lib/types/N-groups.agda},
\code{HoTT-Agda/core/lib/types/Torsor.agda},
\code{HoTT-Agda/theorems/torsors/Delooping.agda} (first 70 lines),
\code{HoTT-Agda/theorems/homotopy/EM-unique.agda},
\code{HoTT-Agda/theorems/homotopy/EM-Grp-equiv.agda} (the equivalence
\code{EM-Ab-equiv}); the statements of \code{orthog-conn-trunc}
(\code{lib/NConnected.agda}), \code{AbGrp-isGrpd}
(\code{lib/groups/Isomorphism.agda}), \code{\ensuremath{\Pi}\ensuremath{\odot}-==-ext}
(\code{lib/types/PtdFibration.agda}), \code{\ensuremath{\odot}EM-Torsors-\ensuremath{\simeq}}
(\code{homotopy/EilenbergMacLane.agda}). Not read: the proofs of
\code{EM-Torsors-\ensuremath{\simeq}}'s dependencies, \code{\ensuremath{\Omega}\textasciicircum{}'S-abgroup}, the EM construction.
Postulates: the repository's own \code{findpostulates} script lists only
HIT and univalence postulates; a full \code{grep postulate} finds, besides lines tagged HIT
or univalence, one untagged block, \code{EilenbergMacLane1/Core.agda} line 44
(\code{emloop-coh'}, \code{EM\ensuremath{_1}-level'}: the coherence and truncation constructors of the
$K(G,1)$ HIT) and \code{Equivalence2.agda} line 223 (``rest of univalence axiom''). The Sính equivalence uses
univalence, function extensionality, truncation and EM HITs; no choice.

\subsection*{Correspondence}

\begin{longtable}{p{0.2\textwidth}p{0.27\textwidth}p{0.47\textwidth}}
\textbf{paper} & \textbf{Agda} & \textbf{statement / Unit counterpart} \\ \hline
Def.~35, Thm~36(1) & \code{Torsors}, \code{tors-ptd-is-trivial} (\code{torsors/Delooping})
 & $T X=\sum_Y \tr Y\times\prod_y X\pteq(Y,y)$; $(\tau=\tau_c)\simeq\mathrm{pr}_1\tau$ when $T_*X$ is contractible.
 Unit: \code{Deloop.Torsors} (small: $Y$ ranges over $\mathrm{BAut}\,X$), \code{Deloop.Torsors.Deloop.F.tr/equiv}
 ($(\mathsf{ptT}=t)\simeq\el\,t$; direction from the centre). \\
Ex.~37 & \code{\ensuremath{\odot}EM-Torsors-\ensuremath{\simeq}}
 & $K(A,n{+}1)\pteq T K(A,n)$. Unit: \emph{by definition}, $K(B,4)=\Tor\,K(B,3)$
 (\code{EM.Step}, \code{Presentation.Pres.K4-def}). \\
Thm~38 & \code{EM-Ab-equiv} (\code{EM-Grp-equiv})
 & $K(-,n):\Ab\simeq\U^*_{\ge n-1,\le n}$, inverse $\pi_n$. Unit: \textbf{new}
 (Thm.~\ref{thm:em}) at $n=2,3$, from \code{Presentation.Rec}, a new \code{Rec2},
 the maps \code{KK.e}, and a new path lemma for \code{Ab}. \\
Cor.~39 & \code{EM-\ensuremath{\Omega}-==-ext} (\code{EM-unique})
 & $(K(A,n)\pteq X)\simeq \text{conn}\times\text{trunc}\times(\pi_nX=A)$. Unit:
 \textbf{not needed} — we work over $\St[1]$ instead of $\Ab$ (Remark~\ref{rem:route}). \\
Lemma~40 & \code{N-Grps-\ensuremath{\simeq}} (\code{lib/types/N-groups}): \code{to}, \code{from}, \code{to-from} (via \code{NG-v2-to-==}), \code{from-to}
 & $\U^*_{\ge0,\le n}\simeq\sum_{B}\sum_{F:B\to\U_{\ge n-1,\le n}}F(\mathrm{pt})$.
 Unit: \textbf{new} at the simply connected level (Prop.~\ref{prop:post}). \\
Thm~41 & \code{NGrp-Sinh-\ensuremath{\simeq}-pre}, \code{NGrp-Sinh-\ensuremath{\simeq}}: \code{orthog-contr}, \code{rearrange},
 \code{sub-rearrange}, \code{EM-\ensuremath{\Omega}-==-ext}, \code{tors-ptd-is-trivial}, \code{\ensuremath{\Pi}\ensuremath{\odot}-==-ext}
 & the chain of equivalences. Unit: \textbf{new} (Prop.~\ref{prop:fib}), simplified:
 the module is constant (simply connected base), the contraction is Lemma~\ref{lem:const}. \\
(in Thm~41) & \code{orthog-conn-trunc}, \code{AbGrp-isGrpd}
 & fillers $X\to\Ab$ over $1$ are contractible for $n$-connected $X$, $\Ab$ a 1-type.
 Unit: Lemma~\ref{lem:const} with $\St[1]$ (a 1-type, Lemma~\ref{lem:stage}) in place of $\Ab$. \\
Note~42 & \code{Sinh-cocycle}: \code{NGrp-Sinh-coc-fib-\ensuremath{\simeq}}
 & $X\simeq\sum_u(\kappa(u)=b_u)$. Unit: the agreement with the realisation
 (Prop.~\ref{prop:real}); same proof (Theorem 36(1) fibrewise). \\
Thm~43 & \code{Sinh-classif-set}
 & set-truncated version. Unit: \textbf{out of scope} (not wanted).
\end{longtable}

\section{Statement}

Unit indices are shifted by two: $\cC\,(m{+}2)$ is $m$-connected, $\cT\,(m{+}2)$ is
$m$-truncated. Throughout, $\ell$ is a level, $\U=\U_\ell$, and
$K(A,m)=\code{KK[A].K}\,m$ with base point $\code{KK[A].pt}\,m$.

\begin{definition}\label{def:types}
\[\SC_\ell = \textstyle\sum_{X:\U_\ell}\sum_{x:X}\cC\,3\,X\times\cT\,5\,X\ :\ \U_{\ell+1},\qquad
\Tri_\ell=\sum_{A:\Ab_\ell}\sum_{B:\Ab_\ell} K(A,2)\to_*K(B,4)\ :\ \U_{\ell+1}.\]
$\Ab_\ell$ is the record \code{Torsor.Ab} (carrier in $\U_\ell$), a type in $\U_{\ell+1}$.
For $n=0,1$, $\St[n]$ is the record \code{EM.Stage[n]}: a type $X:\U_\ell$, a point,
$\cC\,(n{+}3)\,X$, $\cT\,(n{+}4)\,X$ (pointed $(n{+}1)$-connected $(n{+}2)$-types), in $\U_{\ell+1}$.
\end{definition}

The realisation is $\rho(A,B,k)=(\Fib k,(\mathrm{pt},k_0),\code{conn3},\code{trunc5})$ with
$\Fib k=\sum_{x:K(A,2)}k(x)=\mathrm{pt}$ (\code{PostnikovRealization.Triple}).

\begin{theorem}\label{thm:main}
$\rho:\Tri_\ell\to\SC_\ell$ is an equivalence (in $\U_{\ell+1}$).
\end{theorem}

The proof is a chain of equivalences
\[\Tri\ \overset{\text{Thm }\ref{thm:em}}{\simeq}\ \textstyle\sum_{V:\St[0]}\sum_{W:\St[1]}V\to_*\Tor\,W
\ \overset{\text{Prop }\ref{prop:fib}}{\simeq}\ \sum_{V:\St[0]}\Fam\,V
\ \overset{\text{Prop }\ref{prop:post}}{\simeq}\ \SC\]
whose composite $\Phi$ is pointwise equal to $\rho$ (Prop.~\ref{prop:real}); hence $\rho$ is
an equivalence (\code{Equivalences.homotopy} at level $\ell{+}1$).

\section{Generic tools}

\begin{lemma}[paths of pointed types with a property]\label{lem:ptprop}
For a family of propositions $P:\U\to\U$, the type $\sum_{X}\sum_{x:X}P\,X$ has
$((X,x,p)=(Y,y,q))\simeq (X,x)\pteq(Y,y)$, with $\mathsf{idPt}$ over $\refl$. In particular
a pointed equivalence gives a path. Applied to $\SC$ (with $P=\cC\,3\times\cT\,5$) and,
through record $\eta$ (proved by the eliminator \code{split}), to $\St[n]$.
\end{lemma}
\begin{proof} The total space $\sum_{(Y,y,q)}(X,x)\pteq(Y,y)$ is a retract of
$\sum_{(Y,y,e)}P\,Y$ over the contractible \code{PtUniv} (\code{PtEquivOps.pt-univ}), whose fibre
at the centre is the proposition $P\,X$, inhabited by $p$; so it is contractible; the fundamental
theorem of identity types (\code{FundThm.FT}) gives the equivalence. \end{proof}

\begin{lemma}\label{lem:stage} $\St[1]$ is a 1-type: by Lemma~\ref{lem:ptprop} (through the
record) its identity types are equivalent to $(X,x)\pteq(Y,y)$ between 2-connected 3-types, sets
by \code{Lemma4.cor5-equiv} ($t=1$, $c=3$), read at level $\ell{+}1$ by \code{TruncUp.trunc-up}.
\end{lemma}

\begin{lemma}[constant maps]\label{lem:const}
Let $D$ be 1-connected ($\cC\,3$), $Z$ a 1-type ($\cT\,3$), $g:D\to Z$. Then
$\sum_{z:Z}\prod_{d:D}z=g(d)$ is contractible.
\end{lemma}
\begin{proof} Contractibility is a proposition and $D$ is merely inhabited; take $d_0$. The type
is a retract of the pointed sections $\sum_{h:\prod_d g d_0=g d}h(d_0)=\refl$: forth
$(z,h)\mapsto(h',\ \mathsf{trans\text{-}sym\text{-}l})$ with $h'(d)=h(d_0)^{-1}\cdot h(d)$; back
$h\mapsto(g d_0,h)$. The retraction law $(g d_0,h')=(z,h)$ is proved for an arbitrary
$p:z=g d_0$ in place of $h(d_0)$, by $J$ on $p$ (right end fixed, \code{idJR}); at $p=\refl$ it is
$\refl^{-1}\cdot h(d)=h(d)$ pointwise, by function extensionality. The pointed sections are
contractible by \code{Lemma4.lemma4} ($t=0$, $c=2$: a family of sets over a 1-connected base),
at level $\ell{+}1$ since $Z$ is large; its connectedness hypothesis is transported from level
$\ell$ by \code{conn-lift} (induction on the index, comparing the two propositional truncations
with the level-polymorphic eliminator; no resizing). This is Hart--Milner's
\code{orthog-conn-trunc} at the level used. \end{proof}

\begin{lemma}[connectedness of a total space]\label{lem:conn-back}
If $f:A\to B$ has $\cC\,m$ fibres and $B$ is $\cC\,m$, then $A$ is $\cC\,m$.
\end{lemma}
\begin{proof} Induction on $m$, dual to \code{ConnectedMaps.conn-up} (which goes from domain to
codomain): $\tr A$ from $\tr B$ and a fibre;
for $a,a'$, $\ap f:(a=a')\to(fa=fa')$ has $\cC\,(m{-}1)$ fibres
(\code{ConnectedMaps.fib-ap-conn}) and $fa=fa'$ is $\cC\,(m{-}1)$. \end{proof}
For $F:V\to\U$, $\mathrm{pr}_1:\sum_vF v\to V$ has fibres $F v$ (up to the equivalence
$\fib_{\mathrm{pr}_1}(v)\simeq F v$), so $V$ and all $F v$ in $\cC\,3$ give $\sum F$ in $\cC\,3$.

\begin{lemma}[paths in $\Ab$]\label{lem:ab}
An additive bijection $h:G.A\simeq H.A$ gives $G=H$ in $\Ab$.
\end{lemma}
\begin{proof} Induction on the equivalence (\code{Torsor.eqind}) reduces to $h=\mathrm{id}$, with
$G$ and $H$ on the same carrier and $G.\mathrm{add}=H.\mathrm{add}$ pointwise. Then $o$ agree
(\code{EMMaps.hom-unit} of the identity), $\mathrm{inv}$ agree (inverses are unique), $\mathrm{add}$
agree by function extensionality, and the remaining fields are propositions ($\mathrm{setA}$
by \code{Torsor.isProp-isSet}, laws in a set). Record $\eta$ by \code{split} on both sides. \end{proof}

\section{Theorem 38 at $n=2,3$}

Let $\pi_3(W)=\code{Presentation.O3[W.X,W.pt,W.trunc].ab}$ ($\Om^3$ with $\cdot$, commutative by
\code{EH}) and $\pi_2(V)=\code{O2ab}$, the same at $\Om^2V$ (EH at $V$). Let
$K_2(A)=\code{EMn[A].Kstage}\,0$ and $K_3(B)=\code{EMn[B].Kstage}\,1$; by definition
$K_2(A).X=K(A,2)$ and $K_3(B).X=K(B,3)$, and $\Tor\,(K_3B)=\code{Step[1,Kstage 1].next}.X$ is,
by definition, $K(B,4)$ with its base point.

\begin{theorem}\label{thm:em}
$K_2:\Ab\to\St[0]$ and $K_3:\Ab\to\St[1]$ are equivalences, with inverses $\pi_2$, $\pi_3$.
\end{theorem}
\begin{proof} $K_n(\pi_n W)=W$: Lemma~\ref{lem:ptprop} applied to the recognition
equivalence $K(\pi_nW,n)\pteq W$, which is \code{Presentation.Rec.e3} for $n=3$ and a new
\code{Rec2.G2} for $n=2$ (the steps \code{G1}, \code{G2} of \code{Rec}, with $\Om F$ replaced by
$V$; Cor.~12 then Cor.~11 with $c=1$, Lemma 5 = \code{deloop-equiv}).
$\pi_n(K_nB)=B$: Lemma~\ref{lem:ab} applied to the composite
\[\Om^3K(B,3)\xrightarrow{\Om^2 e_2}\Om^2K(B,2)\xrightarrow{\Om e_1}\Om K(B,1)\xrightarrow{e_0}B\]
($e_m=\code{KK.e}\,m$; for $n=2$ the last two), an equivalence (looping of pointed equivalences,
\code{LoopIso.loopPE}) and additive: the addition of $\pi_3=\code{O3.ab}$ is concatenation in
$\Om(\Om^2K(B,3))$, the outermost loop, which $\Om(\Om e_2)$ preserves by
\code{PiOne.Induced.loop-mult}; then $\Om e_1$ the same, and $e_0$ is the band, a homomorphism
(\code{EMSpace.loop0-hom}). \end{proof}
Only the existence of the two paths is used (bi-invertibility), not their values; the
1-truncatedness of $\Ab$ (Hart--Milner's \code{AbGrp-isGrpd}) is \emph{not} used.

\section{The fibrewise step (Theorem 41 at a simply connected base)}

Fix $V:\St[0]$, $X=V.X$, $v_0=V.\mathrm{pt}$. For $W:\St[1]$ let
$\Tor\,W=\sum_{c:\mathrm{BAut}}\tr{\el c}\times\Mu_W(\el c)$ with $\Mu_W\,Y=\prod_{y:Y}(W.X,W.\mathrm{pt})\pteq(Y,y)$
and centre $\mathsf{ptT}_W$ (\code{Deloop.Torsors}, \code{EM.Step[1,W]}). Note the direction: Unit's $\mathsf{tr}_t$ goes
from the centre, $(\mathsf{ptT}=t)\simeq\el\,t.1$, while the paper's Theorem 36(1) reads
$(\tau=\tau_c)\simeq\mathrm{pr}_1\tau$; Prop.~\ref{prop:real} depends on this. Let
\[\textstyle \mathsf{C}\,Y=\cC\,4\,Y\times\cT\,5\,Y,\qquad \Fam\,V=\sum_{F:X\to\U}F(v_0)\times\prod_v\mathsf{C}(F v)\ :\ \U_{\ell+1}.\]

\begin{proposition}\label{prop:fib}
$\sum_{W:\St[1]}(X,v_0)\to_*(\Tor W,\mathsf{ptT}_W)\ \simeq\ \Fam\,V$, by
$(W,k,p)\mapsto(v\mapsto\el(k v).1,\ \mathsf{tr}_{k v_0}(p^{-1}),\ c)$, where
$\mathsf{tr}_t:(\mathsf{ptT}=t)\simeq\el\,t.1$ is Theorem 36(1) (\code{F.tr}, transport of the centre's point)
and $c$ is computed from the structure (step 4 below).
\end{proposition}
\begin{proof} A chain of equivalences.
\begin{enumerate}
\item $(X,v_0)\to_*\Tor W=\sum_{k:X\to\Tor W}k v_0=\mathsf{ptT}\simeq\sum_k\el(k v_0).1$, fibrewise by
$p\mapsto\mathsf{tr}_{k v_0}(p^{-1})$ (\code{sym-equiv} then \code{F.equiv}).
\item $\Tor W\simeq\Tor^L W:=\sum_{Y:\U}\tr Y\times\Mu_W Y$, $(c,h,m)\mapsto(\el c,h,m)$, with
inverse $(Y,h,m)\mapsto(c_Y,\ (h,m)\text{ transported along }p_Y^{-1})$, where
$(c_Y,p_Y:\el c_Y=Y)$ is the centre of $\mathsf{Fib}\,Y=\sum_c\el c=Y$, contractible by
\code{Torsors.fib-contr} from a point of $Y$; since $\mathsf{Fib}\,Y$ is a proposition
(\code{fib-prop}: $\el$ is an embedding) $\tr Y$ suffices (the text of
\code{Presentation.Pres.tor-fib}, re-proved outside \code{Pres}). The two inverse laws: $J$ on
$p_Y$; and $\mathsf{Fib}(\el c)$ being a proposition, $(c_{\el c},p)=(c,\refl)$. Postcomposition gives
$(X\to\Tor W)\simeq(X\to\Tor^LW)$, and \code{SphereAxiom.sigma-base} (a parameter-free lemma of
that file, at level $\ell{+}1$; no sphere axiom) transports the family
$k'\mapsto(k'v_0).1$, which on the image of $k$ is $\el(kv_0).1$ by definition.
\item $\sum_{k':X\to\Tor^LW}(k'v_0).1\ \simeq\ \sum_{F:X\to\U}\sum_{m:\prod_v\tr{Fv}\times\Mu_W(Fv)}F v_0$, by definition
(regrouping; both inverse laws $\refl$).
\item Summing over $W$ and regrouping:
$\sum_W(\ldots)\simeq\sum_F F(v_0)\times S(F)$ with $S(F)=\sum_{W:\St[1]}\prod_v\tr{Fv}\times\Mu_W(Fv)$.
Lemma~\ref{lem:key}: $S(F)\simeq\prod_v\mathsf{C}(Fv)$.
\end{enumerate}\end{proof}

\begin{lemma}[key lemma]\label{lem:key}
For $F:X\to\U$ ($X$ pointed and 1-connected), $S(F)$ and $\prod_v\mathsf{C}(Fv)$ are propositions
and logically equivalent; hence equivalent (\code{prop-equiv}).
\end{lemma}
\begin{proof}
(levels, $S\Rightarrow\mathsf C$) Given $(W,m)$ and $v$: $\mathsf{C}(Fv)$ is a proposition and $\tr{Fv}$ holds; for
$y:Fv$, $m\,v\,y:W.X\simeq Fv$ transfers $\cC\,4$ and $\cT\,5$ from $W$ (\code{to-C}).

($S(F)$ is a proposition) It suffices that $S(F)$ is contractible when inhabited
(\code{prop-of-contr}). From $(W,m)$ we have $c=\code{to-C}(W,m)$, hence
$G:E\to\St[1]$, $G(v,y)=(Fv,y,c\,v)$ with $E=\sum_vFv$. By Lemma~\ref{lem:ptprop}
($\St[1]$ form), $W\pteq G(e)$ is equivalent to $W=G(e)$, so $S(F)$ is a retract of
$\sum_{W}\prod_{e:E}W=G(e)$ (the $\tr{Fv}$ components are propositions, filled by $c$; the
retraction law is \code{sec-eq} of the equivalence, pointwise). The latter is contractible by
Lemma~\ref{lem:const}: $E$ is 1-connected (Lemma~\ref{lem:conn-back}, using $c$) and $\St[1]$
is a 1-type (Lemma~\ref{lem:stage}).

($\mathsf C\Rightarrow S$) Given $c$: $E$ is $\cC\,3$; $Fv_0$ is merely inhabited and $S(F)$ is a
proposition (above), so we may take $y_0:Fv_0$. Put $W_0=(Fv_0,y_0,c\,v_0)$. The family
$Q(v,y)=(Fv_0,y_0)\pteq(Fv,y)$ over $E$ consists of sets (\code{cor5-equiv}), so by
\code{lemma4} ($t=0$, $c=2$) the pointed sections of $Q$ through $\mathsf{idPt}$ at $(v_0,y_0)$
are contractible; the centre gives $\Mu_{W_0}(Fv)$ for all $v$ (the text of
\code{Presentation.Pres.ext}, re-proved outside \code{Pres}).
\end{proof}

\begin{remark}[route]\label{rem:route}
Hart--Milner keep a family $M:X\to\Ab$ and convert $W\pteq(Fv,y)$ into $\pi_3(Fv,y)=M(v)$ by
Corollary 39; the contraction (\code{orthog-contr}) is then fibrewise over $\Ab$. Here the
constancy used is that of $G:E\to\St[1]$ over the 1-connected total space $E$; this one step
covers both their module $M$ and their fibrewise $E_2(u)$, and working with $W:\St[1]$ directly
avoids Corollary 39 and the 1-truncatedness of $\Ab$; $\Ab$ enters only through
Theorem~\ref{thm:em} on the bases. The paper's $T X$ is already large; Unit's $\Tor$ is small
(over $\mathrm{BAut}$), hence step 2.
The base point is carried throughout: the point $F(v_0)$ comes from the pointing $p$ of $k$
through Theorem 36(1); nothing is recovered from connectedness at the end.
\end{remark}

\section{The Postnikov decomposition (Lemma 40)}

Let $\tr{-}_2=\code{(Truncations.TO three)}$ ($\cT\,4$ with $\cC\,4$ fibres of the unit $\eta$).

\begin{proposition}\label{prop:post}
$\alpha:\sum_{V:\St[0]}\Fam\,V\to\SC$, $\alpha(V,F,f,c)=(\sum_vFv,(v_0,f),\cC,\cT)$, is an equivalence,
with inverse $\delta(X,x,c,t)=(V_X,\ \fib_\eta,\ (x,\refl),\ \ldots)$, $V_X=(\tr X_2,\eta x,\ldots)$.
\end{proposition}
\begin{proof}
Levels: $\sum F$ is $\cC\,3$ by Lemma~\ref{lem:conn-back} ($\cC\,4$ fibres lowered by
\code{conn-down}) and $\cT\,5$ by \code{trunc-Sigma} ($V.X$ raised from $\cT\,4$ by \code{trunc-suc});
$\tr X_2$ is $\cC\,3$ by \code{conn-up} along $\eta$ (fibres $\cC\,4$, lowered to $\cC\,3$);
$\fib_\eta(z)$ is $\cC\,4$ (\code{TO.conn}) and $\cT\,5$.

$\alpha\delta=\mathrm{id}$: $w\mapsto w.2.1:\sum_z\fib_\eta z\simeq X$ (the text of
\code{Presentation.Pres.tot-fib}, re-proved generically) sends $(\eta x,(x,\refl))$ to $x$ by
definition; Lemma~\ref{lem:ptprop}.

$\delta\alpha=\mathrm{id}$: let $E=\sum F$, $\hat e:\tr E_2\to V.X$ the extension of
$\mathrm{pr}_1$ (\code{ConnectedMaps.Elim.ext}), $\beta_w:\hat e(\eta w)=w.1$ (\code{ext-beta}).
\begin{itemize}
\item $\hat e$ is an equivalence. Inverse: $g(v)$ is the value of the constant map
$y\mapsto\eta(v,y):Fv\to\tr E_2$ (\code{ConnectedMaps.null}: $Fv$ is $\cC\,4$ and $\tr E_2$ is
$\cT\,4$, so constants are an equivalence), with $g_\beta:g(v)=\eta(v,y)$. $\hat e g=\mathrm{id}$:
at $v$, the map $y\mapsto\ap_{\hat e}(g_\beta)\cdot\beta_{(v,y)}:Fv\to(\hat e(gv)=v)$ into a
$\cT\,3\subseteq\cT\,4$ type is constant (\code{null} again). $g\hat e=\mathrm{id}$:
\code{Elim.ext} along $\eta$ into $z\mapsto g(\hat ez)=z$, given at $\eta w$ by
$\ap_g(\beta_w)\cdot g_\beta$.
\item $t_z:\fib_\eta(z)\simeq F(\hat e z)$ is the composite
$\fib_\eta(z)\simeq\sum_{w}\hat e(\eta w)=\hat ez$ ($\ap_{\hat e}$, an equivalence by
\code{ap-equiv}) $\simeq\sum_w w.1=\hat ez$ (precomposition with $\beta_w^{-1}$)
$\simeq F(\hat ez)$ (\code{fib-pr1}), so that, by definition,
\[t_z(w,q)=\operatorname{transport}_F\bigl(\beta_w^{-1}\cdot\ap_{\hat e}(q)\bigr)(w.2).\]
\item Coherence, with $w_0=(v_0,f)$: $t_{\eta w_0}(w_0,\refl)=\operatorname{transport}_F(\beta_{w_0}^{-1})(f)$
by definition (\code{t-pt}), and transporting this along the pointing $\beta_{w_0}:\hat e(\eta w_0)=v_0$
of $\hat e$ gives $f$ (\code{subst-cancel}).
\end{itemize}
These data give the path by an SIP for $\sum_{V}\Fam V$ (\code{sip}): induction on the pointed
equivalence (\code{StageId.J}, from the contractibility of $\sum_T S\pteq T$) reduces to the
identity, where $F'=F$ by function extensionality and univalence from $t$ (\code{fam-P}), the
point by \code{ua-beta} and the coherence, and $\prod_v\mathsf C$ is a proposition (\code{fam-path}).
\end{proof}

\section{Agreement with the realisation (Note 42)}

\begin{proposition}\label{prop:real}
For $(A,B,k):\Tri$, $\Phi(A,B,k)=\rho(A,B,k)$ in $\SC$.
\end{proposition}
\begin{proof} By the definitions of the steps, $\Phi(A,B,k)=(\sum_{x:K(A,2)}\el(kx).1,\ (\mathrm{pt},\mathsf{tr}_{k\,\mathrm{pt}}(k_0^{-1})),\ldots)$
(Theorem~\ref{thm:em} enters only by $A\mapsto K_2A$, $B\mapsto K_3B$, definitionally the same
spaces). The map $(x,q)\mapsto(x,\mathsf{tr}_{kx}(q^{-1})):\Fib k\to\sum_x\el(kx).1$ is an
equivalence (\code{Fibres.sigma-equiv} of \code{sym-equiv} and \code{F.equiv}, as in
\code{Presentation.Pres.tot}) and sends $(\mathrm{pt},k_0)$ to the base point by definition
(\code{Fib} uses $k(x)=\mathrm{pt}$; hence the inverse $q^{-1}$). Lemma~\ref{lem:ptprop}. This is Hart--Milner's
\code{NGrp-Sinh-coc-fib-\ensuremath{\simeq}}. \end{proof}

\section{Consequence: $\pi_3(S^2)=\mathbb{Z}$ through the classification}

\code{MapFib.Final.rep-fib} (under the $S^2$ axiom) proves the marked universal property
$s^*:[P,F]_*\simeq[S^2,F]_*$, $P=\Fib\kappa$, for every fibre $F=\Fib k$ of a triple. By
Theorem~\ref{thm:main}, every $X:\SC$ is $\rho(\mathsf{sec}\,X)$ up to the path
$\mathsf{sec\text{-}eq}\,X:\rho(\mathsf{sec}\,X)=X$ in $\SC$; transporting \code{rep-fib} at the
triple $\mathsf{sec}\,X$ along it gives $s^*:[P,X]_*\simeq[S^2,X]_*$ for every pointed simply
connected 3-type $X$ (\code{Pi3Classified.Final.rep}). This is the interface
\code{Pi3Sphere.Conditional.MapThm}, and \code{Pi3Sphere.Conditional.pi3-S2-T} (with the constructed
3-truncation of $S^2$) gives \code{Pi3Classified.Final.pi3-S2}$\,:\pi_3(S^2)\cong\mathbb{Z}$ as groups,
from the $S^2$ axiom alone, without the special presentation of \code{Presentation.Pres}. No
classification of maps is needed: the mapping theorem supplies that part.

\section{Declarations (checked)}

All in \code{sag/pi3-s2-postnikov}, inside the segment \code{Theory} (parameters: univalence,
images, surjectivity of images).

\begin{longtable}{p{0.42\textwidth}p{0.52\textwidth}}
\code{ClassAb}: \code{R}, \code{Props}, \code{rest-prop}, \code{mk}, \code{ab-eta}, \code{mk-eq}, \code{ab-path} & Lemma~\ref{lem:ab} \\
\code{ClassBase.PtProp[P].I.equiv/path}, \code{StageId[n].I.equiv/path}, \code{StageId.J} & Lemma~\ref{lem:ptprop} \\
\code{ClassBase.stage-gpd} & Lemma~\ref{lem:stage} \\
\code{ClassBase.conn-lift}, \code{Const}, \code{const-contr} & Lemma~\ref{lem:const} \\
\code{ClassBase.conn-back}, \code{fib-pr1}, \code{total-conn} & Lemma~\ref{lem:conn-back} \\
\code{ClassBase.raise}, \code{prop-of-contr}, \code{prop-equiv}, \code{pair-path}, \code{ap-equiv}, \code{subst-cancel} & generic \\
\code{ClassEM.O2ab}, \code{Rec2}, \code{K2}, \code{K3}, \code{pi2}, \code{pi3}, \code{k2-pi2}, \code{k3-pi3},
\code{Band2}, \code{Band3} (\code{fun-def}, \code{add-hom}), \code{pi2-k2}, \code{pi3-k3}, \code{k2-equiv}, \code{k3-equiv} & Theorem~\ref{thm:em} \\
\code{ClassFib.C}, \code{Fam}, \code{TorW[W]} (\code{tau}, \code{toL}, \code{ofL}, \code{toL-equiv}) & Prop.~\ref{prop:fib}, steps 1--2 \\
\code{ClassFib.Key[X,v0,cX,F]}: \code{to-C}, \code{Uniq.contr}, \code{s-prop}, \code{Exist}, \code{of-C}, \code{key} & Lemma~\ref{lem:key} \\
\code{ClassFib.Chain[V]}: \code{e01}, \code{e12}, \code{e23}, \code{e34}, \code{fib-equiv} & Prop.~\ref{prop:fib} \\
\code{ClassPost.SC3}, \code{alpha}, \code{Dl.delta}, \code{alpha-delta}, \code{fam-path}, \code{sip},
\code{DA[u]} (\code{ehat}, \code{ehat-equiv}, \code{t}, \code{t-pt}, \code{coh}), \code{delta-alpha}, \code{post-equiv} & Def.~\ref{def:types}, Prop.~\ref{prop:post} \\
\code{PostnikovClassification.Triples}, \code{realise}, \code{base}, \code{chain} & Def.~\ref{def:types}, the chain \\
\code{PostnikovClassification.Agree}, \code{realise-agree} & Prop.~\ref{prop:real} \\
\code{PostnikovClassification.realise-equiv}, \code{classification} & Theorem~\ref{thm:main} \\
\code{Pi3Classified.Final}: \code{rep}, \code{map-thm}, \code{pi3-S2} & \S9 \\
\end{longtable}

The final statements are
\begin{verbatim}
realise-equiv  : P1.isEquiv Triples CP.SC3 realise
classification : P1.Equiv Triples CP.SC3 = (realise, realise-equiv)
\end{verbatim}
in \code{U (l+1)}, with $\SC$ the type \code{ClassPost.SC3} (no alias is declared, rule 2). The
inverse and both inverse laws are the fields \code{sec}, \code{sec-eq}, \code{ret},
\code{ret-eq} of \code{realise-equiv}.

\section{Scope, checks and residual diagnostics}
The equivalence transports identity types, hence equivalences of the classified objects; it is
not a classification of noninvertible maps nor an equivalence of homotopy categories.
Assumptions: only the Theory parameters; no sphere axiom, choice, resizing or new HIT. The new
modules use parameter-free top-level declarations of sphere-named files (\code{SphereAxiom.O2},
\code{sigma-base}; \code{SphereUniversal.three}, \code{five}) and of \code{Presentation}
(\code{O3}, \code{Rec}, \code{deloop-equiv}); none of the segments \code{SphereAxiom.Sphere},
\code{Presentation.Pres/Final}, \code{Sphere3*}, \code{MapFib}, \code{CupSquare} is used by the
classification. Only \code{Pi3Classified} (\S9) uses \code{MapFib.Final} and \code{CupSquare.kappa},
inside its segment \code{Final} whose parameters are the $S^2$ axiom.

Checks (28 October 2026): fresh \code{seg All.seg}: 4461 declarations, 0 failed, 2.23\,s;
server \code{segc sag/pi3-s2-postnikov all}: 89 modules, 0 failed, 1.96\,s; uniform checker
\code{segd --opens}: all 20 specifications (the 19 existing ones and \code{classification.opens},
355 entries, including \code{Pi3Classified}) with no deviation. Every new declaration checks in at most 0.01\,s.
\code{classification.opens} was generated from \code{--opened} and read against this note module by
module; it is not a line-by-line textual audit. One residual diagnostic: the \code{refl} lemma
\code{ClassEM.Band2.fun-def} (``$\mathsf{fun}\,x=\mathsf{loop0}(\Om e_1\,x)$ by definition'') opens
the band's internals (\code{EM1.band}, \code{TrivG.kb/tb}, \code{TorId.AutM.fa},
\code{Trunc.induction}) when comparing \code{(KK.e 0).1.1} with \code{EMSpace.loop0.1}; it is
isolated in that one declaration, and \code{Band2.add-hom}, \code{Band3.add-hom} no longer open it.
\end{document}
